% ------------------------------------------------------------------
% Binary NORM-versus-MI Classification -- 10-seed Results subsection
% (values computed from 10 trained seeds per model; see Table~\ref{tab:binary_summary})
% ------------------------------------------------------------------

\subsection{Binary NORM-versus-MI Classification}
\label{sec:binary_results}

The explainability experiments presented in the following sections are based on a binary
NORM-versus-MI classification task. This simplified setting facilitates the interpretation
of attribution maps and enables a more direct comparison of the different attribution
methods. Classification performance was evaluated on the held-out test set over ten
independent random initialisations, using Accuracy, Macro F1, ROC-AUC, and Average
Precision. Table~\ref{tab:binary_summary} summarises the mean performance across all ten
seeds, and Figure~\ref{fig:binary_boxplots} shows the corresponding variability.

The results in Table~\ref{tab:binary_summary} indicate that the two architectures achieved
highly comparable performance on the binary task. The CNN--BiLSTM with Attention obtained
marginally higher mean values on all four metrics, namely Accuracy (0.912 vs.\ 0.907),
Macro F1 (0.904 vs.\ 0.897), ROC-AUC (0.972 vs.\ 0.969), and Average Precision
(0.960 vs.\ 0.956). These differences are small and of the same order as the seed-to-seed
variation, so they should not be over-interpreted; overall, both models achieved comparable
predictive performance on the held-out test set.

The boxplots in Figure~\ref{fig:binary_boxplots} confirm that the two metric distributions
overlap substantially for every metric, further indicating that the observed differences are
minor. The variability across seeds was small for both architectures. ROC-AUC was the most
stable metric, with standard deviations of approximately 0.001--0.002, suggesting consistent
ranking performance across different random initialisations, whereas Accuracy and Macro F1
exhibited slightly greater dispersion. In all cases the distributions of the two models
remain largely overlapping.

Figure~\ref{fig:binary_roc_pr} shows the mean ROC and Precision--Recall curves (with MI as
the positive class) aggregated over the ten seeds, with $\pm$one standard-deviation bands.
The ROC curves of the two architectures almost completely overlap across the full operating
range, and the Precision--Recall curves differ only marginally at high recall. Both models
therefore exhibit almost identical discrimination performance on the binary NORM-versus-MI
task, consistent with the near-identical ROC-AUC and Average Precision values reported in
Table~\ref{tab:binary_summary}.

% ---- Captions (for the two figures) ----
% Figure (binary boxplots):
% "Distribution of Accuracy, Macro F1 score, ROC-AUC, and Average Precision for the CNN
%  and CNN-BiLSTM with Attention models across ten random seeds. Each box shows the
%  interquartile range and median; the ten individual seed values are overlaid as jittered
%  points."
%
% Figure (binary ROC/PR):
% "Macro-average one-versus-rest ROC and Precision-Recall curves for the CNN and
%  CNN-BiLSTM with Attention across ten random seeds."
